\section{Consequences for cut selection}\label{disc:section}

The unrestricted corner bound is a useful reference value. It says what
one quadratic constraint can enforce on a chosen basis cone, and it
separates a loss caused by the relaxation from a loss caused by the free
set family. For low-dimensional constraints, the support and inertia
results make this reference value accessible through small subproblems.
For larger projected dimension, the complexity result explains why one
should not expect a general efficient exact selection procedure.

The determinant orbit provides a small optimization problem and a
geometric explanation for its failures. Its parameters can be searched
by semidefinite feasibility, its contacts can force a one-dimensional
pencil, and exact certificates can prove that every parameter misses a
target bound. This makes the family useful even when it is not exact:
it supplies a controlled baseline against which a point rule, a
completion, or a larger family can be assessed. The minor examples show
that this role extends beyond a single bilinear lifting variable.

The scaled-depth guarantee is conditional in an explicit and verifiable
way. If $D$ is bounded across a collection of bilinear corners, the orbit
retains a positive fraction of each corner bound. If $D$ grows, the sharp
examples show that a fixed fraction cannot be promised. The additional
conditioning term for a fixed point rule explains a further loss that
can occur without changing the bilinear constraint. Translation and
rescaling can alter that rule, so its guarantee should be applied to the
actual quadratic representation used in the implementation.

Completion and closure address different weaknesses. Upward completion
can remove artificial finite exits and can make a family exact at a
corner where the uncompleted orbit has an infinite approximation
factor. A closure can combine cuts with complementary coefficient
patterns. Neither operation supplies general exactness: the certified
closure examples retain a gap, and the sharp family shows why there is
no uniform guarantee independent of scaled vertex depth. A better
single-cut value and a better closure value should therefore be measured
separately.

For repeated cuts, local objective improvement is insufficient. The
convergence theorem requires a positive depth whenever violation stays
away from zero. The explicit stalling sequence retains fixed violation
and well-behaved basis cones while its depths vanish. It demonstrates a
failure of arbitrary maximal-set selection. It does not establish that
SCIP's default rule stalls, nor that the implemented approximate orbit
search satisfies the theorem's uniform-depth hypotheses. These are
distinct assertions requiring distinct evidence.

The numerical comparisons reinforce this separation. Stronger corner
bounds can lead to weaker full-LP bounds after reoptimization, and small
changes in successive corners can change later progress. Reduced-cost
floors used in a practical search also replace the exact objective
criterion by a surrogate. The retained solver studies provide no
consistent improvement from the tested corner-bound and efficacy
searches. The full-solve records contain implementation and timing
qualifications that prevent a general recommendation about enabling
intersection cuts. They remain useful evidence about the tested rules
and about the measurements needed in future implementations.

Several extensions would require additional results rather than a
stronger reading of the present ones. The support-one criterion for
the uncompleted orbit does not characterize exactness of every upward
completion. The depth guarantee for $\famB$ does not automatically give
the same guarantee for its transformed point-rule subfamily. The
three-parameter determinant orbit does not describe every maximal
free set of signature $(2,2)$. Finally, the idealized convergence rule
does not certify an approximate implementation without a uniform
approximation and cone-geometry bound. The results proved here delimit
these questions and provide exact examples and quantitative tests for
addressing them.
